% Living Showcase: semantische Tabellen und kompakte Grids (Palette Slot 4).
\StartChapter[ch:tables]{Tabellen und Grids}

\SubsectionBar[sec:table-header]{Header und Zebra}

\begin{tablebox}[Farbiger Header, Zebra ab Zeile 2\ZSFTitleTag{tablebox + ZSFtable}]
  \begin{ZSFtable}{Y{0.98} Y{1.32} Q{0.70}}
    \ZSFheaderRow \ZSFhead{Element} & \ZSFhead{Wirkung} & \ZSFhead{Zeile} \\
    \ZSFkeyword{ZSFtable} & Header und Zebra & 1 \\
    ZSFheaderRow & färbt die Kopfzeile & 1 \\
    ZSFhead & färbt jede Kopfzelle & 1 \\
    Zebra & ab der zweiten Zeile & 2
  \end{ZSFtable}
\end{tablebox}

\SubsectionBar[sec:table-variants]{Ohne Header oder Zebra}

\begin{tablebox}[Zebra ab der ersten Zeile\ZSFTitleTag{ZSFtableFlat}]
  \begin{ZSFtableFlat}{Y{1.2} Q{0.8}}
    \ZSFkeyword{ZSFtableFlat} & ohne Kopfzeile \\
    Streifen & ab Zeile 1 \\
    Einsatz & reine Datenlisten
  \end{ZSFtableFlat}
\end{tablebox}

\begin{tablebox}[Ganz ohne Streifen\ZSFTitleTag{ZSFtablePlain}]
  \begin{ZSFtablePlain}{Y{1.25} Z{0.75}}
    \ZSFkeyword{ZSFtablePlain} & kein Zebra \\
    Einsatz & sehr kurze Tabellen
  \end{ZSFtablePlain}
\end{tablebox}

\begin{tablebox}[Spaltentypen im Vergleich\ZSFTitleTag{L/C/R/F}]
  \begin{ZSFtablePlain}{L C R F{1}}
    L linksbündig & C zentriert & R rechtsbündig & F{1}: x^2+1
  \end{ZSFtablePlain}
\end{tablebox}

\begin{valuegrid}{2}[Kompaktes Wertegitter\ZSFTitleTag{valuegrid}]
  \ZSFkeyword{valuegrid} & Spalte 1 & Spalte 2 \\
  Argument & Anzahl & Datenspalten \\
  Erste Spalte & immer & Beschriftung
\end{valuegrid}

\begin{runintext}
  Alle Spaltenbreiten sind proportional: die Faktoren einer Tabelle müssen
  zur Anzahl X-Spalten summieren, dann füllt sie die Spalte exakt aus.
\end{runintext}
